% Claim proof file for row claim_3_10 -- "TwoDisjointNode".
% Auto-generated stub by scaffold/solve_chapter.py at row start. A manager
% agent dedicated to writing the TeX proof fills in the proof body below.
%
% This file is a sibling of `main.tex` inside `<subsection>/tex/`. The
% `subfiles` package lets it render both standalone and as part of
% `main.tex` via `\subfile{...}`.
%
% The statement is restated above the proof so the file is self-contained
% when rendered on its own. The proof must:
%   - Stay close to the lecture notes' proof if one exists (always search
%     the LN first for a `\begin{proof}` block immediately following the
%     claim; copy / lightly edit when present).
%   - Otherwise use the LN paradigm and cite prior claims by their `ref`.
%   - Be self-contained enough that a mathematician can follow it without
%     re-reading the LN.
%   - Have no `sorry`, `omitted`, or "obvious" shortcuts.

\documentclass[main]{subfiles}

\begin{document}
% ref: claim_3_10 (proof, with statement restated for readability)
% title: TwoDisjointNode
\def\rowref{claim\_3\_10}
\phantomsection\label{claim_3_10}

% --- statement (restated) ---------------------------------------------------
\begin{Lem}[Two disjoint node-splitting hard interventions commute]
   Let $G=(J,V,E,L)$ be a CADMG and $W_1, W_2 \ins V$ two disjoint subsets of the output nodes from $G$. 
   Then the CADMG obtained from first node-splitting on $W_1$ and then node-splitting on $W_2$ is the same CADMG that arises from first node-splitting on $W_2$ and then node-splitting on $W_1$:
   \[ \lp G_{\swig(W_1)} \rp_{\swig(W_2)} =  \lp G_{\swig(W_2)} \rp_{\swig(W_1)}  =  G_{\swig(W_1 \dcup W_2)}.   \]
\end{Lem}

% --- proof ------------------------------------------------------------------
\begin{proof}
% TODO: write the proof body.
\end{proof}

\end{document}
